% ============================================================
%  The Particle Model — Lesson Booklet (one section per lesson: A Reading · B Theory · C Practice · D Above level)
%  Year 7 (Year 8 extension)  |  Victorian Curriculum F–10 Version 2.0
%  Compile with: xelatex three times, then python3 fill.py Particle_Model_Lesson_Booklet_Yr7
% ============================================================
\documentclass[11pt,a5paper,twoside]{article}
\newcommand{\unitname}{Advanced Science Skills}\newcommand{\unitcaps}{ADVANCED SCIENCE SKILLS}

% ── Fonts (compile with LuaLaTeX) ────────────────────────────────────────────
\usepackage{fontspec}
% Bundled fonts (in ./fonts/) so the document compiles anywhere with XeLaTeX/LuaLaTeX
\setmainfont{Lato}[                   % body — humanist sans (stands in for deck's Inter)
  Path=../fonts/, Extension=.ttf, Ligatures=TeX,
  UprightFont=Lato-Regular, ItalicFont=Lato-Italic,
  BoldFont=Lato-Bold, BoldItalicFont=Lato-BoldItalic ]
\newfontfamily\headingfont{Poppins}[  % display/headings — matches the slide decks
  Path=../fonts/, Extension=.ttf, Ligatures=TeX,
  UprightFont=Poppins-Regular, ItalicFont=Poppins-Italic,
  BoldFont=Poppins-Bold, BoldItalicFont=Poppins-BoldItalic ]
\setlength{\headheight}{22pt}

\usepackage[protrusion=true,expansion=false]{microtype}
\usepackage[a5paper, inner=17mm, outer=13mm, top=19mm, bottom=17mm, headheight=14pt, headsep=4mm, footskip=8mm]{geometry}
\usepackage{xcolor}
\usepackage{tikz}
\usetikzlibrary{arrows.meta, shapes.geometric, calc, positioning, decorations.pathreplacing}
\usepackage{tcolorbox}
\tcbuselibrary{skins, breakable, raster}
\usepackage{fancyhdr}
\usepackage{titlesec}
\usepackage{enumitem}
\usepackage{booktabs}
\usepackage{tabularx}
\usepackage{array}
\usepackage{colortbl}
\usepackage{multicol}
\usepackage{amsmath}
\usepackage{needspace}
% TikZ diagram macros are embedded below (just before \begin{document}) for easy editing.
\usepackage{longtable}
\usepackage{graphicx}
\usepackage{pifont}
\usepackage{lastpage}
\usepackage{calc}
\usepackage{mdframed}
\usepackage{setspace}
\usepackage{ragged2e}
% PDF bookmarks: one per lesson (made in \sectiondiv) plus the cover and
% contents, and links from the contents table to each lesson.
% implicit=false: hyperref adds no anchors of its own (an anchor on each
% list item nudged lists down 1.5 mm, which would undo the page fill).
% hidelinks keeps the links invisible, so the printed pages are unchanged.
\usepackage[implicit=false, hidelinks,
  pdftitle={Advanced Science Skills},
  pdfsubject={Years 7–8 Science, Victorian Curriculum F–10 v2.0}]{hyperref}
\usepackage{bookmark} % writes the outline itself, which implicit=false switches off in hyperref
% hyperref 7.01 with implicit=false leaves \Hy@PageAnchorSlide undefined, which
% throws at \end{document} and stops fill.py (it builds with -halt-on-error).
\makeatletter\providecommand\Hy@PageAnchorSlide{}\makeatother

% ── Colour Palette ───────────────────────────────────────────────────────────
% Remapped to the slide-deck palette (names kept so body + diagrams cascade)
\definecolor{NavyBlue}   {HTML}{16305C}
\definecolor{MidBlue}    {HTML}{2E6FE0}
\definecolor{LightBlue}  {HTML}{E8F1FF}
\definecolor{TealGreen}  {HTML}{0F6E63}
\definecolor{LightGreen} {HTML}{E0F7F3}
\definecolor{GoldDark}   {HTML}{C77F12}
\definecolor{GoldLight}  {HTML}{FFF3DC}
\definecolor{PurpleDark} {HTML}{5B3FC0}
\definecolor{PurpleLight}{HTML}{EFE9FF}
\definecolor{OrangeDark} {HTML}{C9483D}
\definecolor{OrangeLight}{HTML}{FFE9E6}
\definecolor{MidGrey}    {HTML}{AEBED4}
\definecolor{LightGrey}  {HTML}{F2F5FA}
\definecolor{DarkGrey}   {HTML}{3A4759}
\definecolor{TableHead}  {HTML}{16305C}
\definecolor{RowAlt}     {HTML}{EEF4FF}
% Bright accent tints (deck) for bars, rules and chips
\definecolor{Blue}   {HTML}{2E6FE0}
\definecolor{Teal}   {HTML}{18A999}
\definecolor{Amber}  {HTML}{F5A623}
\definecolor{Coral}  {HTML}{FF6B5E}
\definecolor{Violet} {HTML}{7C5CE0}
\definecolor{Sky}    {HTML}{4DA3FF}
\definecolor{Ink}    {HTML}{16223A}
\definecolor{Line}   {HTML}{DDE5F0}
\definecolor{Paper}  {HTML}{FBFCFE}

% ── Page Style ───────────────────────────────────────────────────────────────
\pagestyle{fancy}
\fancyhf{}
\renewcommand{\headrulewidth}{0.4pt}
\renewcommand{\headrule}{\hbox to\headwidth{\color{NavyBlue}\leaders\hrule height \headrulewidth\hfill}}
\fancyhead[L]{\headingfont\small\color{NavyBlue}The Particle Model}
\fancyhead[R]{\headingfont\small\color{DarkGrey}Year 7 \textbar{} Victorian Curriculum F--10 v2.0}
\fancyfoot[C]{\headingfont\small\color{DarkGrey}\thepage\ \textbar{} \pageref{LastPage}}
\fancypagestyle{plain}{%
  \fancyhf{}
  \fancyfoot[C]{\headingfont\small\color{DarkGrey}\thepage\ \textbar{} \pageref{LastPage}}
  \renewcommand{\headrulewidth}{0pt}
}

% ── Section Formatting ───────────────────────────────────────────────────────
\titleformat{\section}
  {\headingfont\color{NavyBlue}\Large\bfseries}
  {}{0em}{}
  [{\color{Teal}\rule[0.6ex]{2.2em}{3pt}}{\color{Line}\rule[0.6ex]{\dimexpr\linewidth-2.2em\relax}{1pt}}]
\titlespacing{\section}{0pt}{18pt}{8pt}

\titleformat{\subsection}
  {\headingfont\color{MidBlue}\large\bfseries}
  {\thesubsection}{0.6em}{}
\titlespacing{\subsection}{0pt}{12pt}{4pt}

\titleformat{\subsubsection}
  {\headingfont\color{NavyBlue}\normalsize\bfseries}
  {}{0em}{}
\titlespacing{\subsubsection}{0pt}{8pt}{2pt}

% Deck-style eyebrow / kicker (optional use)
\newcommand{\kicker}[2][Teal]{%
  {\headingfont\bfseries\footnotesize\color{#1}%
   \textcolor{#1}{\rule[0.18ex]{1.6em}{2.2pt}}\hspace{0.5em}\MakeUppercase{#2}\par}%
  \vspace{2pt}}

% ── Custom tcolorbox Environments ────────────────────────────────────────────

% Section divider banner
\newtcolorbox{sectionbanner}[2]{
  enhanced, breakable=false,
  colback=NavyBlue, colframe=NavyBlue,
  fontupper=\color{white}\bfseries\large,
  left=4mm, right=4mm, top=3mm, bottom=3mm,
  title={\color{white}\Huge\bfseries #1},
  attach boxed title to top left={yshift=-2mm, xshift=4mm},
  boxed title style={colback=NavyBlue, colframe=NavyBlue},
  after upper={\par\smallskip\normalfont\normalsize\color{white!80!NavyBlue}\textit{#2}}
}

% Key info box (blue)
\newtcolorbox{infobox}[1]{
  breakable, arc=2.5mm, boxrule=0.5pt,
  colback=LightBlue, colframe=Blue!40,
  fonttitle=\headingfont\bfseries, coltitle=MidBlue,
  title={#1},
  left=4mm, right=3.5mm, top=2.5mm, bottom=2.5mm
}

% Scaffold box (yellow)
\newtcolorbox{scaffold}[1]{
  breakable, arc=2.5mm, boxrule=0.5pt,
  colback=GoldLight, colframe=Amber!40,
  fonttitle=\headingfont\bfseries, coltitle=GoldDark,
  title={#1},
  left=4mm, right=3.5mm, top=2.5mm, bottom=2.5mm
}

% Worked example box
\newtcolorbox{workedexbox}[1]{
  breakable, arc=2.5mm, boxrule=0.5pt,
  colback=white, colframe=Amber!50,
  fonttitle=\headingfont\bfseries, coltitle=GoldDark,
  title={#1},
  left=4mm, right=3.5mm, top=2.5mm, bottom=2.5mm
}
% Worked examples: \begin{workedex}{title} -> "Worked Example N.M --- title",
% auto-numbered per section and labelled we:N.M. Never type the number.
\newcounter{wenum}[section]
\renewcommand{\thewenum}{\thesection.\arabic{wenum}}
\newenvironment{workedex}[1]
  {\refstepcounter{wenum}\label{we:\thewenum}\begin{workedexbox}{Worked Example \thewenum\ --- #1}}
  {\end{workedexbox}}

% We do: straight after each worked example, a near-identical question the
% class works through together, with its own writing space. Numbered to
% match the worked example above it and labelled wd:N.M.
%   \begin{wedo}{title} question ... \answerbox{h} \end{wedo}
\newtcolorbox{wedobox}[1]{
  enhanced, breakable=false, arc=2.5mm, boxrule=0.5pt,
  colback=white, colframe=TealGreen!55,
  fonttitle=\headingfont\bfseries, coltitle=TealGreen!75!black,
  title={#1},
  left=4mm, right=3.5mm, top=2.5mm, bottom=2.5mm
}
\newenvironment{wedo}[1]
  {\par\needspace{9\baselineskip}\begin{wedobox}{We do \thewenum\ --- #1}\label{wd:\thewenum}\ignorespaces}
  {\end{wedobox}}

% Year 7 box (blue)
\newtcolorbox{yr7box}{
  breakable, arc=2.5mm, boxrule=0.5pt,
  colback=LightBlue, colframe=Blue!40,
  left=4mm, right=3.5mm, top=2.5mm, bottom=2.5mm
}

% Year 8 extension box (violet)
\newtcolorbox{yr8box}{
  breakable, arc=2.5mm, boxrule=0.5pt,
  colback=PurpleLight, colframe=Violet!40,
  left=4mm, right=3.5mm, top=2.5mm, bottom=2.5mm
}

% Answer lines box
\newcommand{\answerbox}[1]{%
  \begin{tcolorbox}[
    enhanced, breakable, arc=2mm,
    colback=Paper, colframe=Line,
    left=2mm, right=2mm, top=1mm, bottom=1mm,
    boxrule=1pt,
    height=#1
  ]
  \end{tcolorbox}%
}

% Experiment header
\newtcolorbox{expheader}[3]{
  breakable=false, arc=2.5mm,
  colback=RowAlt, colframe=Amber!50, boxrule=0.5pt,
  fonttitle=\headingfont\bfseries\large, coltitle=NavyBlue,
  title={Experiment #1: #2},
  left=4mm, right=3.5mm, top=2.5mm, bottom=2.5mm,
  after upper={\par\smallskip
    {\headingfont\bfseries\color{GoldDark}Aim:}\ \textit{#3}},
}

% Reading comprehension passage box (green)
\newtcolorbox{readingbox}[1]{
  breakable, arc=2.5mm, boxrule=0.5pt,
  colback=white, colframe=Teal!50,
  fonttitle=\headingfont\bfseries, coltitle=TealGreen,
  title={#1},
  left=4mm, right=3.5mm, top=2.5mm, bottom=2.5mm,
  fontupper=\normalsize
}

% Skill-builder box (for beginner scaffolding)
\newtcolorbox{skillbox}[1]{
  breakable, arc=2.5mm, boxrule=0.5pt,
  colback=RowAlt, colframe=Blue!40,
  fonttitle=\headingfont\bfseries, coltitle=MidBlue,
  title={#1},
  left=4mm, right=3.5mm, top=2.5mm, bottom=2.5mm
}

% ── Custom Commands ───────────────────────────────────────────────────────────

% Question: \question{marks}{text}  (marks can be empty)
% Numbered automatically as Q<section>.<n>, reset every section.
% Do NOT write numbers by hand.
\newcounter{qnum}[section]
\renewcommand{\theqnum}{\thesection.\arabic{qnum}}
\newcommand{\question}[2]{%
  \refstepcounter{qnum}%
  \needspace{3\baselineskip}%
  \medskip
  \noindent\label{q:\theqnum}\textbf{Q\theqnum}\ifx&#1&\else\ \textcolor{DarkGrey}{(#1 mark\ifnum#1=1\else s\fi)}\fi\quad #2%
  \medskip
}

% Inline marks note
\newcommand{\markscount}[1]{\hfill\textcolor{DarkGrey}{\textit{(#1 mark\ifnum#1=1\else s\fi)}}}

% ── Practice set counter and question macros ─────────────────────────────────
\newcounter{pqnum}
\newcommand{\practiceset}[1]{%
  \needspace{6\baselineskip}
  \setcounter{pqnum}{0}%
  \vspace{6pt}
  {\headingfont\bfseries\large\color{GoldDark} Practice Set --- #1}\\[-2pt]
  {\color{Amber}\rule{\linewidth}{1pt}}\\[2pt]
}
% \pq{question text}{working box height}
\newcommand{\pq}[2]{%
  \stepcounter{pqnum}%
  \needspace{4\baselineskip}
  \vspace{5pt}
  \noindent\textbf{\thepqnum.}\quad #1\par
  \vspace{2pt}
  \answerbox{#2}
}
% Two short questions side by side, each with its own box
\newcommand{\pqpair}[4]{%
  \needspace{4\baselineskip}
  \vspace{5pt}
  \noindent
  \begin{minipage}[t]{0.48\linewidth}
    \stepcounter{pqnum}\textbf{\thepqnum.}\quad #1\par\vspace{2pt}
    \begin{tcolorbox}[enhanced,arc=2mm,colback=Paper,colframe=Line,left=2mm,right=2mm,
      top=1mm,bottom=1mm,boxrule=1pt,height=#2]\end{tcolorbox}
  \end{minipage}\hfill
  \begin{minipage}[t]{0.48\linewidth}
    \stepcounter{pqnum}\textbf{\thepqnum.}\quad #3\par\vspace{2pt}
    \begin{tcolorbox}[enhanced,arc=2mm,colback=Paper,colframe=Line,left=2mm,right=2mm,
      top=1mm,bottom=1mm,boxrule=1pt,height=#4]\end{tcolorbox}
  \end{minipage}
}

% ── Comprehension question macro (compact, lined) ────────────────────────────
\newcounter{rcnum}
\newcommand{\rcreset}{\setcounter{rcnum}{0}}
\newcommand{\rcq}[2]{%
  \stepcounter{rcnum}%
  \needspace{3\baselineskip}
  \vspace{4pt}
  \noindent\textbf{C\thercnum.}\quad #1\par
  \vspace{2pt}
  \linedbox{#2}
}

% Section divider (big numbered banner): \sectiondiv{title}{summary}
%
% This IS the section: it steps the real section counter, so subsections,
% tasks, questions and figures inside it are numbered N.M automatically.
% Put \label{sec:N} straight after it. Never hand-type a section number.
\newcommand{\sectiondiv}[2]{%
  \needspace{5\baselineskip}%
  \refstepcounter{section}%
  % PDF bookmark for the lesson; its anchor, lesson.N.1, is what the contents table links to
  \pdfbookmark[1]{Lesson \thesection{} — #1}{lesson.\thesection}%
  \begin{tcolorbox}[
    enhanced, breakable=false, arc=3.5mm,
    colback=NavyBlue, colframe=NavyBlue, boxrule=0pt,
    sidebyside, sidebyside align=center,
    lefthand width=2.4cm,
    left=2mm, right=4mm, top=4mm, bottom=4mm
  ]
    \begin{center}
      {\headingfont\bfseries\color{white!35!NavyBlue}\fontsize{50}{50}\selectfont \thesection}
    \end{center}
    \tcblower
    {\headingfont\bfseries\color{white}\LARGE #1}\\[3pt]
    {\itshape\color{white!75!NavyBlue}\small #2}
  \end{tcolorbox}%
}

% Tasks: \task, \task[a] (part a of a NEW task), \taskpart{b} (next part
% of the same task), \exttask (Year 8 "Extension" — same counter).
% Numbered Task <section>.<n>, reset every section. Do NOT type numbers.
% Every task, question and figure also gets an automatic label
% (task:1.7b, q:2.4, fig:1.2, fig:E2.1) so page refs come from the .aux.
\newcounter{tasknum}[section]
\renewcommand{\thetasknum}{\thesection.\arabic{tasknum}}
\NewDocumentCommand{\task}{o}{\refstepcounter{tasknum}\label{task:\thetasknum\IfValueT{#1}{#1}}\textbf{Task \thetasknum\IfValueT{#1}{#1}}}
\newcommand{\taskpart}[1]{\label{task:\thetasknum#1}\textbf{Task \thetasknum#1}}
\newcommand{\exttask}{\refstepcounter{tasknum}\label{task:\thetasknum}\textbf{Extension \thetasknum}}

% Figure captions: \figcap[fig:label]{caption} -> "Figure N.M --- caption"
% \expfigcap{caption} -> "Figure EN.M" (experiment pages). Both auto-numbered.
\newcounter{fignum}[section]
\renewcommand{\thefignum}{\thesection.\arabic{fignum}}
\newcounter{expfignum}[section]
\renewcommand{\theexpfignum}{E\thesection.\arabic{expfignum}}
\newcommand{\figcap}[2][]{%
  \begin{center}
    \refstepcounter{fignum}\label{fig:\thefignum}\if\relax\detokenize{#1}\relax\else\label{#1}\fi
    \small\itshape\color{DarkGrey} Figure \thefignum\ --- #2
  \end{center}%
}
\newcommand{\expfigcap}[1]{%
  \begin{center}
    \refstepcounter{expfignum}\label{fig:\theexpfignum}%
    \small\itshape\color{DarkGrey} Figure \theexpfignum\ --- #1
  \end{center}%
}

% Diagram include
\newcommand{\diagramfig}[3]{%  {filename}{width fraction}{caption}
  \begin{center}
    \includegraphics[width=#2\linewidth]{#1}\\[4pt]
    \small\itshape\color{DarkGrey} #3
  \end{center}%
}

% Table column types
\newcolumntype{L}[1]{>{\RaggedRight\arraybackslash}p{#1}}
\newcolumntype{C}[1]{>{\Centering\arraybackslash}p{#1}}

% Table header colour
\newcommand{\thead}[1]{\cellcolor{TableHead}\color{white}\bfseries #1}
\newcommand{\talt}{\rowcolor{RowAlt}}
\newcommand{\tmid}{\rowcolor{MidGrey}}

% Inline year tags
\newcommand{\tagyr}[1]{%
  \tikz[baseline=-0.5ex]\node[
    rounded corners=2pt,
    fill=#1!20, draw=#1, text=#1,
    inner xsep=4pt, inner ysep=1pt,
    font=\scriptsize\bfseries
  ]{Yr #1 only};}

% Draw answer lines
\newcommand{\anslines}[1]{%
  \vspace{2pt}%
  \foreach \n in {1,...,#1}{%
    \par\noindent\rule{\linewidth}{0.3pt}\vspace{3pt}%
  }%
}

% ── Document ─────────────────────────────────────────────────────────────────
% ── Optional-lesson note, shown under the lesson banner ──────────────────────
\newcommand{\optionalnote}[1]{%
  % the title is a bold lead-in on the first line rather than a title bar,
  % which keeps Lesson 11's reading on the lesson's first page
  \begin{tcolorbox}[enhanced, breakable=false, arc=2.5mm, boxrule=0.5pt,
    colback=PurpleLight, colframe=Violet!50,
    left=4mm, right=3.5mm, top=2mm, bottom=2mm]
  {\headingfont\bfseries\color{PurpleDark}Optional lesson \textbullet{} Year 8 extension.}\enspace #1
  \end{tcolorbox}\vspace{2pt}}


% ══════════════════════════════════════════════════════════════════════════════
%  LESSON STRUCTURE (lesson booklet, Oct 2026)
%  Every lesson runs in the same four parts, in this order:
%    \lessonpart{A}{Reading and comprehension}{Teal}   \label{read:N}
%    \lessonpart{B}{Theory and examples}{Blue}         \label{theory:N}
%    \lessonpart{C}{Practice}{Amber}                    \label{prac:N}
%    \lessonpart{D}{Above level}{Violet}                \label{ext:N}
%  Above-level / higher-order questions live ONLY in Part D, at the end.
% ══════════════════════════════════════════════════════════════════════════════
\newcommand{\lessonpart}[3]{%
  \par\if#1A\needspace{10\baselineskip}\vspace{10pt}\else
    \if#1C\notesfill\fi\clearpage\fi
  \begin{tcolorbox}[enhanced, breakable=false, arc=2mm, boxrule=0pt,
    colback=#3!14, borderline west={4pt}{0pt}{#3},
    left=4mm, right=3mm, top=1.6mm, bottom=1.6mm]
  {\headingfont\bfseries\color{#3!60!black}\large Part #1 \textbullet{} #2}
  \end{tcolorbox}\vspace{2pt}}

% Key vocabulary at the start of each lesson: students copy each definition
% from the vocabulary slide in the lesson's starter. First row = red scaffold.
\newcommand{\vocabhead}{%
  \par\vspace{6pt}%
  \begin{tcolorbox}[enhanced, breakable=false, arc=2mm, boxrule=0pt,
    colback=NavyBlue!10, borderline west={4pt}{0pt}{NavyBlue},
    left=4mm, right=3mm, top=1.6mm, bottom=1.6mm]
  {\headingfont\bfseries\color{NavyBlue}\large Key vocabulary}
  \end{tcolorbox}\vspace{2pt}%
  \noindent Copy the definition of each word from the slide. The first one has been done for you.\par}

% Curriculum codes under a lesson banner
\newcommand{\lessoncodes}[1]{%
  {\small\headingfont\color{DarkGrey}Victorian Curriculum 2.0:\ \textbf{#1}}\par\vspace{4pt}}

% Two short tasks side by side, each with its own working box
% \taskpair{text}{box height}{text}{box height}
% The two boxes line up: the shorter question is padded by the difference in
% text height (measured in a scratch box, task counter restored, no label).
\newsavebox\tpbox
\newlength\tpA \newlength\tpB
\newcount\tpsaved
\newcommand\tpmeasure[2]{%
  \tpsaved=\value{tasknum}%
  \sbox\tpbox{\vbox{\hsize=0.485\linewidth \linewidth=\hsize \parskip=0pt
    \renewcommand\label[1]{}\task #1\par}}%
  \setlength{#2}{\dimexpr\ht\tpbox+\dp\tpbox\relax}%
  \setcounter{tasknum}{\tpsaved}}
\newcommand{\taskpair}[4]{%
  \par\needspace{9\baselineskip}\vspace{4pt}%
  \tpmeasure{#1}{\tpA}\tpmeasure{#3}{\tpB}%
  \noindent
  \edef\abgrp{\number\numexpr\value{abox}+1\relax}%
  \begin{minipage}[t]{0.485\linewidth}
    \task #1\par\ifdim\tpB>\tpA\vspace{\dimexpr\tpB-\tpA\relax}\fi\vspace{2pt}\answerbox{#2}
  \end{minipage}\hfill
  \begin{minipage}[t]{0.485\linewidth}
    \task #3\par\ifdim\tpA>\tpB\vspace{\dimexpr\tpA-\tpB\relax}\fi\vspace{2pt}\answerbox{#4}
  \end{minipage}\par\gdef\abgrp{0}\vspace{4pt}}
% Two parts of one question side by side, with their answer boxes level:
% \ptpair{a}{part (a) text}{box height}{b}{part (b) text}{box height}.
% The shorter part's text gets space below it to match the longer one.
\newcommand\ptmeasure[3]{%
  \sbox\tpbox{\vbox{\hsize=0.485\linewidth \linewidth=\hsize \parskip=0pt \pt{#1}#2\par}}%
  \setlength{#3}{\dimexpr\ht\tpbox+\dp\tpbox\relax}}
\newcommand{\ptpair}[6]{%
  \ptmeasure{#1}{#2}{\tpA}\ptmeasure{#4}{#5}{\tpB}%
  \par\noindent
  \edef\abgrp{\number\numexpr\value{abox}+1\relax}%
  \begin{minipage}[t]{0.485\linewidth}
    \pt{#1}#2\par\ifdim\tpB>\tpA\vspace{\dimexpr\tpB-\tpA\relax}\fi\answerbox{#3}
  \end{minipage}\hfill
  \begin{minipage}[t]{0.485\linewidth}
    \pt{#4}#5\par\ifdim\tpA>\tpB\vspace{\dimexpr\tpA-\tpB\relax}\fi\answerbox{#6}
  \end{minipage}\par\gdef\abgrp{0}}

% Answer boxes never split across a page (a half box on its own page is useless)
% ── Page filling ────────────────────────────────────────────────────────────
% Every writing box gets an id. Its page is logged to \jobname.abp; fill.py
% measures the empty space left at the foot of each page and shares it out
% among that page's boxes via \jobname.abx (\abset{id}{extra}). Side-by-side
% boxes share a group so they grow together.
\newcounter{abox}
\newwrite\abpos \immediate\openout\abpos=\jobname.abp
\newcommand{\abset}[2]{\expandafter\gdef\csname abx@#1\endcsname{#2}}
\InputIfFileExists{\jobname.abx}{}{}
\def\abgrp{0}
\newcommand{\abnext}{\stepcounter{abox}%
  \edef\abcur{\ifcsname abx@\theabox\endcsname\csname abx@\theabox\endcsname\else 0pt\fi}}
\newcommand{\ablog}{\edef\abtmp{\noexpand\write\abpos{\theabox\space\abgrp\space\noexpand\thepage}}\abtmp}
\newcommand{\abh}[1]{\dimexpr#1+\abcur\relax}
\renewcommand{\answerbox}[1]{%
  \abnext
  \begin{tcolorbox}[enhanced, breakable=false, arc=2mm, colback=Paper, colframe=Line,
    left=2mm, right=2mm, top=1mm, bottom=1mm, boxrule=1pt, height=\abh{#1}]
  \end{tcolorbox}\ablog%
}
% \linedbox{height}: an answer box ruled in dotted thirds. Lines are \thirdrow/3 apart and run dotted, dotted,
% solid (the baseline), repeating, and end on one dotted line. The number of
% rows is worked out from the box's final height, so the ruling still fills the
% box after fill.py has grown it; the rows are centred in the box.
\newlength\thirdrow \setlength\thirdrow{9mm}
\tikzset{thirdsolid/.style={draw=DarkGrey!70, line width=0.5pt},
         thirddot/.style={draw=DarkGrey!70, line width=0.6pt, dash pattern=on 0.9pt off 1.6pt}}
\newlength\lbH \newlength\lbO \newcount\lbN
\newcommand{\linedbox}[1]{%
  \abnext
  % interior height = box height less the 1pt frame top and bottom
  \setlength\lbH{\dimexpr\abh{#1}-2pt\relax}%
  % whole rows that fit with at least 3 mm spare top and bottom (at least one)
  \lbN=\numexpr(\dimexpr\lbH-6mm-\thirdrow/2\relax)/\thirdrow\relax
  \ifnum\lbN<1 \lbN=1 \fi
  \setlength\lbO{\dimexpr(\lbH-\lbN\thirdrow)/2\relax}%
  \begin{tcolorbox}[enhanced, breakable=false, arc=2mm, colback=Paper, colframe=Line,
    left=2mm, right=2mm, top=1mm, bottom=1mm, boxrule=1pt, height=\abh{#1},
    underlay={\begin{tcbclipinterior}
      \foreach \i in {1,...,\the\lbN}{
        \draw[thirddot] ([xshift=3mm,yshift=-\lbO-(\i-1)*\thirdrow]interior.north west) -- ([xshift=-3mm,yshift=-\lbO-(\i-1)*\thirdrow]interior.north east);
        \draw[thirddot] ([xshift=3mm,yshift=-\lbO-(\i-1)*\thirdrow-\thirdrow/3]interior.north west) -- ([xshift=-3mm,yshift=-\lbO-(\i-1)*\thirdrow-\thirdrow/3]interior.north east);
        \draw[thirdsolid] ([xshift=3mm,yshift=-\lbO-(\i-1)*\thirdrow-2*\thirdrow/3]interior.north west) -- ([xshift=-3mm,yshift=-\lbO-(\i-1)*\thirdrow-2*\thirdrow/3]interior.north east);}
      \draw[thirddot] ([xshift=3mm,yshift=-\lbO-\lbN\thirdrow]interior.north west) -- ([xshift=-3mm,yshift=-\lbO-\lbN\thirdrow]interior.north east);
    \end{tcbclipinterior}}]
  \end{tcolorbox}\ablog%
}
% Every written-response box is ruled the same way (Part A comprehension, We do,
% Part C tasks and Part D extensions, including side-by-side pairs). Space for
% mathematical working uses \workbox{height} instead and stays plain, as do
% \drawbox and the write-in tables.
\let\workbox\answerbox   % plain box: mathematical working (calculations) stays unruled
\let\answerbox\linedbox
% \drawbox{height}{contents}: white box for a diagram or drawing
\newcommand{\drawbox}[2]{%
  \abnext
  \begin{tcolorbox}[colback=white, colframe=MidGrey, boxrule=0.5pt, height=\abh{#1}]#2\end{tcolorbox}\ablog}
% Notes box at the end of Part B: fills whatever is left of the page (if worth it)
\newcommand{\notesfill}{\par
  \ifdim\dimexpr\pagegoal-\pagetotal\relax>3.2cm
    \vspace*{\fill}%
    \begin{tcolorbox}[enhanced, breakable=false, arc=2mm, colback=Paper, colframe=Line,
      left=2mm, right=2mm, top=1mm, bottom=1mm, boxrule=1pt,
      height=\dimexpr\pagegoal-\pagetotal-1.0cm\relax]
      {\small\headingfont\color{DarkGrey}Notes}
    \end{tcolorbox}%
  \fi}



% Centred table: wrap every content table in \begin{ctab} ... \end{ctab}
% (figures already centre themselves inside their \diag... macros)
% ctab also registers write-in tables (S column) for page filling: the extra
% height is per row, logged with the row count.
\newenvironment{ctab}{\par\smallskip\begingroup\centering
  \abnext\let\tbcur\abcur\global\wtrows0 }{\par
  \ifnum\wtrows>0 \edef\abtmp{\noexpand\write\abpos{\theabox\space\abgrp\space\noexpand\thepage\space\the\wtrows}}\abtmp\fi
  \endgroup\smallskip}

% ── Booklet formatting rules (Oct 2026) ─────────────────────────────────────
% 1. No marks are shown anywhere in the booklet.
\renewcommand{\markscount}[1]{}
% 2. A task's heading sits on its own line; its text starts on the next
%    line, left aligned. Parts (a), (b)… each start a new line: use \pt{a}.
\RenewDocumentCommand{\task}{o}{\par\noindent\refstepcounter{tasknum}\label{task:\thetasknum\IfValueT{#1}{#1}}\textbf{Task \thetasknum\IfValueT{#1}{#1}}\par\nobreak\noindent\ignorespaces}
\renewcommand{\taskpart}[1]{\par\noindent\label{task:\thetasknum#1}\textbf{Task \thetasknum#1}\par\nobreak\noindent\ignorespaces}
\renewcommand{\exttask}{\par\noindent\refstepcounter{tasknum}\label{task:\thetasknum}\textbf{Extension \thetasknum}\par\nobreak\noindent\ignorespaces}
\newcommand{\pt}[1]{\par\noindent\textbf{(#1)}\ \ignorespaces}
% Workstation pages (Experiment 1): numbered method steps with a hanging
% indent, and a one-line safety reminder under the materials list.
\newcommand{\wsstep}[1]{\par\noindent\hangindent=1.6em\hangafter=1\makebox[1.6em][l]{\textbf{#1}}\ignorespaces}
\newcommand{\wssafety}[1]{\par\smallskip{\small\textcolor{ScaffoldRed}{\headingfont\bfseries Safety:}\ #1}\par}
% 3. Tables have vertical rules between columns and down both sides
%    (write | between and around the columns, e.g. {|L{3cm}|C{2cm}|}).
%    Keep the normal padding at the outer edges (no @{} there): colortbl paints
%    header colour one \tabcolsep past each edge, so trimming the padding makes
%    the header overhang the grid.
\arrayrulecolor{MidGrey}
\setlength{\arrayrulewidth}{0.6pt}
% 4. Write-in tables: W/Z are left/centred columns that centre text
%    vertically; put the S column first (see below).
\newcolumntype{W}[1]{>{\RaggedRight\arraybackslash}m{#1}}
\newcolumntype{Z}[1]{>{\Centering\arraybackslash}m{#1}}
\def\tbcur{0pt}
\newcount\wtrows
\newcommand{\wh}{\global\advance\wtrows1 \rule[\dimexpr-0.75cm-\tbcur/2\relax]{0pt}{\dimexpr1.5cm+\tbcur\relax}}
% S = zero-width first column that makes every row of a write-in table
% at least 1.5 cm tall with its text vertically centred. Start each data
% row with "&"; start the header row with "\hd &" so it stays slim.
\newif\ifnowh
\newcolumntype{S}{@{}>{}c<{\ifnowh\else\wh\fi\global\nowhfalse}@{\hspace{\tabcolsep}}}
\newcommand{\hd}{\global\nowhtrue}
% 5. The first row of a multi-row write-in table is a worked scaffold in red.
\definecolor{ScaffoldRed}{HTML}{C0392B}
\newcommand{\scaf}[1]{\textcolor{ScaffoldRed}{\itshape #1}}
% 6. Tables are a light grid (rules between rows and columns) on white —
%    row shading is off, because colortbl paints shading over column rules.
\renewcommand{\talt}{}
\renewcommand{\tmid}{}
% ══════════════════════════════════════════════════════════════════════════════
%  DIAGRAM DEFINITIONS — The Particle Model
%  Self-contained TikZ macros, no external image files. Each \diag... macro is
%  called from the body by name.
% ══════════════════════════════════════════════════════════════════════════════
\usetikzlibrary{arrows.meta,calc,positioning,patterns,backgrounds}
\tikzset{
  thin arrow/.style={-{Stealth[length=6pt,width=4pt]}, line width=1.4pt},
}

% ── Lesson 1: the same three jars holding a solid, a liquid and a gas ───────────
% The solid keeps its own shape, the liquid takes the jar's shape with a flat top,
% the gas fills the whole jar. No particles yet: those arrive in Lesson 4.
\newcommand{\jar}[1]{% jar outline with its base at (#1,0)
  \draw[line width=1.2pt, Ink] (#1-1.25,2.9) -- (#1-1.25,0.15) arc (180:270:0.15)
    -- (#1+1.1,0) arc (270:360:0.15) -- (#1+1.25,2.9);}
\newcommand{\diagStates}{%
\begin{center}
\begin{tikzpicture}[font=\small, x=1cm, y=1cm]
  % solid: a block sitting in the corner
  \fill[GoldLight] (-0.95,0.02) rectangle (0.15,1.25);
  \draw[GoldDark, line width=1.1pt] (-0.95,0.02) rectangle (0.15,1.25);
  \jar{0}
  \node[font=\bfseries, text=NavyBlue] at (0,-0.45) {Solid};
  \node[align=center, font=\footnotesize, text=DarkGrey] at (0,-1.05) {keeps its own shape\\and volume};
  % liquid: takes the jar's shape, flat top
  \begin{scope}[xshift=4.2cm]
    \fill[Sky!35] (-1.25,0.15) arc (180:270:0.15) -- (1.1,0) arc (270:360:0.15) -- (1.25,1.1) -- (-1.25,1.1) -- cycle;
    \draw[MidBlue, line width=1.1pt] (-1.25,1.1) -- (1.25,1.1);
    \jar{0}
    \node[font=\bfseries, text=NavyBlue] at (0,-0.45) {Liquid};
    \node[align=center, font=\footnotesize, text=DarkGrey] at (0,-1.05) {takes the jar's shape,\\keeps its volume};
  \end{scope}
  % gas: fills the whole (sealed) jar
  \begin{scope}[xshift=8.4cm]
    \fill[Violet!16] (-1.25,0.15) arc (180:270:0.15) -- (1.1,0) arc (270:360:0.15) -- (1.25,2.9) -- (-1.25,2.9) -- cycle;
    \draw[line width=2.2pt, DarkGrey] (-1.4,2.9) -- (1.4,2.9);
    \jar{0}
    \node[font=\scriptsize, text=DarkGrey, anchor=south] at (0,2.95) {lid};
    \node[font=\bfseries, text=NavyBlue] at (0,-0.45) {Gas};
    \node[align=center, font=\footnotesize, text=DarkGrey] at (0,-1.05) {spreads out to fill\\the whole jar};
  \end{scope}
\end{tikzpicture}
\end{center}}

% ── Lesson 1, Worked Example 1.2 and We do 1.2: the same syringe, before and after
\newcommand{\diagSyringe}[3]{% {start mL}{end mL}{contents: gas|liquid}
\begin{center}\resizebox{\linewidth}{!}{%
\begin{tikzpicture}[font=\scriptsize, x=0.11cm, y=1cm]
  \foreach \row/\vol/\lab in {0/#1/before, -1.15/#2/after}{
    \begin{scope}[yshift=\row cm]
      \ifx#3g\fill[Violet!16] (0,0) rectangle (\vol,0.7);\else\fill[Sky!35] (0,0) rectangle (\vol,0.7);\fi
      \draw[line width=1pt, Ink] (0,0) rectangle (60,0.7);
      \draw[line width=1pt, Ink] (0,0.25) -- (-4,0.25) (0,0.45) -- (-4,0.45);
      \fill[Ink] (-5.5,0.15) rectangle (-4,0.55);
      \foreach \t in {0,10,...,60}{\draw[Ink] (\t,0.7) -- (\t,0.85); \node[anchor=south, font=\tiny] at (\t,0.82) {\t};}
      \fill[DarkGrey] (\vol,0.04) rectangle (\vol+1.6,0.66);
      \draw[line width=2pt, DarkGrey] (\vol+1.6,0.35) -- (\vol+26,0.35);
      \draw[line width=2pt, DarkGrey] (\vol+26,0.05) -- (\vol+26,0.65);
      \node[anchor=east, font=\scriptsize\bfseries, text=DarkGrey] at (-7,0.35) {\lab};
    \end{scope}}
  \node[font=\tiny, text=DarkGrey, anchor=west] at (61,0.35) {mL};
  \node[font=\tiny, text=DarkGrey, anchor=north] at (-4.75,-1.2) {sealed tip};
\end{tikzpicture}}
\end{center}}

% needspace.sty reserves room by planting a penalty -100 with 9 lines of
% temporary stretch. On a page that can't be filled tightly that fake break
% always wins, so a We do that fitted was pushed over and left half a page
% empty. This version breaks only when the room really isn't there, and
% ragged bottom lets every other page run on until something won't fit;
% fill.py then grows the writing boxes into what is left at the foot.
\renewcommand{\needspace}[1]{\par\ifdim\dimexpr\pagegoal-\pagetotal\relax<\dimexpr#1\relax\newpage\fi}
\raggedbottom
% Part D sits in a breakable yr8box, where \needspace can't see the page (the box
% is built off the page first). Wrap each extension in keepext so it is one
% unsplittable block: the yr8box then breaks between extensions, never inside one.
\newenvironment{keepext}{\par\begin{tcolorbox}[blankest, breakable=false]}{\end{tcolorbox}\par}
\usetikzlibrary{patterns}
\providecommand{\lessonart}{\thesection}

\providecommand{\unitname}{The Particle Model}\providecommand{\unitcaps}{THE PARTICLE MODEL}
% ══════════════════════════════════════════════════════════════════════════════
%  2027 WORKBOOK STYLE — A5 TOPIC BOOKLET (added by make_a5.py)
%  Same look as the Space A5 trial and the A4 experiment booklets: colour
%  cover, black-and-white inside, Outfit headings, a lesson opener band with
%  the lesson's ink symbol, boxes told apart by frame and icon:
%    I do ● solid · We do ●● dashed · You do ●●● heavy · key idea / reading ◆
%    thin · Above level ▲ double · Country ◎ dotted
%  The body still names the master's colours; they are mapped to ink here.
% ══════════════════════════════════════════════════════════════════════════════
\usepackage{pdflscape}
\usetikzlibrary{patterns}

% A picture wider than the A5 text block is scaled down to fit it.
\newsavebox\fitb
\newenvironment{fitbox}{\begin{lrbox}{\fitb}}{\end{lrbox}%
  \ifdim\wd\fitb>\linewidth\resizebox{\linewidth}{!}{\usebox\fitb}\else\usebox\fitb\fi}

% ── Fonts: Outfit for headings ──────────────────────────────────────────────
\renewfontfamily\headingfont{Outfit}[
  Path=../fonts/, Extension=.ttf, Ligatures=TeX,
  UprightFont=Outfit-Regular, BoldFont=Outfit-SemiBold,
  ItalicFont=Outfit-Regular, ItalicFeatures={FakeSlant=0.18},
  BoldItalicFont=Outfit-SemiBold, BoldItalicFeatures={FakeSlant=0.18} ]

% ── Black and white inside ──────────────────────────────────────────────────
\definecolor{Ink}{HTML}{111111}
\colorlet{NavyBlue}{Ink}\colorlet{MidBlue}{Ink}\colorlet{TealGreen}{Ink}
\colorlet{GoldDark}{Ink}\colorlet{PurpleDark}{Ink}\colorlet{OrangeDark}{Ink}
\colorlet{TableHead}{Ink}
\colorlet{Blue}{Ink}\colorlet{Teal}{Ink}\colorlet{Amber}{Ink}\colorlet{Coral}{Ink}
\colorlet{Violet}{Ink}\colorlet{Sky}{Ink}
\colorlet{LightBlue}{white}\colorlet{LightGreen}{white}\colorlet{GoldLight}{white}
\colorlet{PurpleLight}{white}\colorlet{OrangeLight}{white}\colorlet{RowAlt}{white}
\colorlet{LightGrey}{white}\colorlet{Paper}{white}
\colorlet{DarkGrey}{black!72}\colorlet{MidGrey}{black!45}\colorlet{Line}{black!45}
\arrayrulecolor{black!60}
% Particle kinds stay distinguishable without colour: white, light grey,
% dotted and dark.
\tikzset{
  pA/.style={circle, fill=white, draw=Ink, line width=0.6pt, minimum size=3.2mm, inner sep=0pt},
  pB/.style={circle, fill=black!35, draw=Ink, line width=0.6pt, minimum size=3.2mm, inner sep=0pt},
  pC/.style={circle, preaction={fill=white}, pattern=dots, draw=Ink, line width=0.6pt, minimum size=3.2mm, inner sep=0pt},
  pD/.style={circle, fill=black!75, draw=Ink, line width=0.6pt, minimum size=3.2mm, inner sep=0pt}}

% ── Page style: mirrored, running head on the outside ───────────────────────
\fancypagestyle{fancy}{\fancyhf{}%
  \renewcommand{\headrulewidth}{0.4pt}%
  \renewcommand{\headrule}{\hbox to\headwidth{\color{Ink}\leaders\hrule height \headrulewidth\hfill}}%
  \fancyhead[LE]{\headingfont\footnotesize \unitname}%
  \fancyhead[RO]{\headingfont\footnotesize\leftmark}%
  \fancyfoot[LE,RO]{\headingfont\footnotesize\thepage}}
\fancypagestyle{plain}{\fancyhf{}\renewcommand{\headrulewidth}{0pt}%
  \fancyfoot[LE,RO]{\headingfont\footnotesize\thepage}}
\fancypagestyle{opener}{\fancyhf{}\renewcommand{\headrulewidth}{0pt}%
  \fancyfoot[LE,RO]{\headingfont\footnotesize\thepage}}
\pagestyle{fancy}\markboth{}{}
\raggedbottom
\setlength{\emergencystretch}{2em}

\titleformat{\subsection}{\headingfont\large\bfseries}{\thesubsection}{0.6em}{}
\titleformat{\subsubsection}{\headingfont\normalsize\bfseries}{}{0em}{}

% ── Lesson opener band (replaces the navy section banner) ────────────────────
\renewcommand{\sectiondiv}[2]{%
  \clearpage\thispagestyle{opener}%
  \refstepcounter{section}\markboth{\thesection\ \textbullet\ #1}{\thesection\ \textbullet\ #1}%
  \pdfbookmark[1]{Lesson \thesection{} — #1}{lesson.\thesection}\hypertarget{lesson.\thesection.1}{}%
  \noindent\begin{minipage}[b]{\dimexpr\linewidth-36mm\relax}
    {\headingfont\bfseries\scriptsize\addfontfeatures{LetterSpace=14}\unitcaps\ \textbullet\ TOPIC}\par\vspace{1mm}
    {\headingfont\bfseries\fontsize{40}{40}\selectfont\thesection}\par\vspace{1.5mm}
    {\headingfont\bfseries\fontsize{15}{18}\selectfont\RaggedRight #1\par}\vspace{1mm}
    {\footnotesize\itshape\RaggedRight #2\par}
  \end{minipage}\hfill
  \begin{minipage}[b]{34mm}\includegraphics[width=34mm]{../art/l\lessonart-ink.pdf}\end{minipage}\par
  \vspace{2.5mm}\noindent\rule{\linewidth}{1.6pt}\par\vspace{1mm}}
\renewcommand{\lessoncodes}[1]{{\small\itshape VC 2.0: #1}\par\vspace{2pt}}

% Lesson parts: a heading and a rule; B, C and D start a new page.
\renewcommand{\lessonpart}[3]{%
  \par\if#1A\needspace{10\baselineskip}\vspace{8pt}\else
    \if#1C\notesfill\fi\clearpage\fi
  {\headingfont\bfseries\Large Part #1 \textbullet\ #2}\par\vspace{0.5mm}\noindent\rule{\linewidth}{0.6pt}\par\vspace{1mm}}
\renewcommand{\vocabhead}{%
  \par\vspace{4pt}{\headingfont\bfseries\Large Key vocabulary}\par\vspace{1pt}%
  \noindent{\small Copy the meaning of each word from the slide. The first one has been done for you.}\par}
\renewcommand{\thead}[1]{\headingfont\bfseries #1}

% ── Box system ──────────────────────────────────────────────────────────────
\newcommand{\icodots}[1]{\tikz[baseline=-0.6ex]{\foreach \k in {1,...,#1}{\fill[white] (\k*1.9mm,0) circle (0.7mm);}}}
\newcommand{\icokey}{\tikz[baseline=-0.6ex]{\fill[white] (0,0.9mm)--(0.9mm,0)--(0,-0.9mm)--(-0.9mm,0)--cycle;}}
\newcommand{\icoup}{\tikz[baseline=-0.6ex]{\fill[white] (-1mm,-0.8mm)--(1mm,-0.8mm)--(0,1mm)--cycle;}}
\newcommand{\icocountry}{\tikz[baseline=-0.6ex]{\draw[white,line width=0.5pt] (0,0) circle (1mm); \fill[white] (0,0) circle (0.45mm);}}
\newcommand{\icoread}{\tikz[baseline=-0.6ex]{\draw[white,line width=0.6pt] (-1.1mm,-0.8mm) rectangle (1.1mm,0.8mm); \draw[white,line width=0.5pt] (0,-0.8mm) -- (0,0.8mm);}}
\newcommand{\tagicon}[1]{\tikz[baseline=-0.6ex]\node[fill=Ink, rounded corners=0.8mm, inner xsep=1.2mm, inner ysep=1mm]{#1};}
\tcbset{wbbase/.style={enhanced, colback=white, arc=1.2mm,
    left=3mm, right=3mm, top=4.2mm, bottom=2.4mm, before skip=10pt, after skip=7pt,
    attach boxed title to top left={xshift=3mm, yshift=-2.4mm},
    boxed title style={colback=Ink, colframe=Ink, arc=0.8mm, boxrule=0pt,
                       left=1.6mm, right=1.8mm, top=0.6mm, bottom=0.6mm},
    fonttitle=\headingfont\bfseries\footnotesize, coltitle=white}}
\renewtcolorbox{workedexbox}[1]{wbbase, breakable, colframe=Ink, boxrule=1pt,
  title={\icodots{1}\ I do \textbullet\ #1}}
\renewenvironment{workedex}[1]
  {\refstepcounter{wenum}\label{we:\thewenum}\begin{workedexbox}{Worked example \thewenum\ --- #1}}
  {\end{workedexbox}}
\renewtcolorbox{wedobox}[1]{wbbase, breakable=false, frame hidden,
  borderline={1.1pt}{0pt}{Ink, dash pattern=on 4pt off 2.6pt},
  title={\icodots{2}\ #1}}
\renewenvironment{wedo}[1]
  {\par\begin{wedobox}{We do \thewenum\ --- #1}\label{wd:\thewenum}\ignorespaces}
  {\end{wedobox}}
\renewtcolorbox{yr7box}{wbbase, breakable, colframe=Ink, boxrule=2.4pt, top=5mm,
  title={\icodots{3}\ You do \textbullet\ Check your understanding}}
\renewtcolorbox{yr8box}{wbbase, breakable, frame hidden,
  borderline={0.6pt}{0pt}{Ink}, borderline={0.6pt}{1.6pt}{Ink},
  title={\icoup\ Above level \textbullet\ Year 8 extension}}
\renewtcolorbox{infobox}[1]{wbbase, breakable, colframe=Ink, boxrule=0.5pt, title={\icokey\ #1}}
\renewtcolorbox{scaffold}[1]{wbbase, breakable, colframe=Ink, boxrule=0.5pt, title={\icokey\ #1}}
\renewtcolorbox{skillbox}[1]{wbbase, breakable, colframe=Ink, boxrule=0.5pt, title={\icokey\ #1}}
\renewtcolorbox{readingbox}[1]{wbbase, breakable, colframe=Ink, boxrule=0.5pt, title={\icoread\ #1}}
\renewtcolorbox{expheader}[3]{wbbase, breakable=false, colframe=Ink, boxrule=1pt,
  title={\icokey\ Experiment #1: #2},
  after upper={\par\smallskip{\headingfont\bfseries Aim:}\ \textit{#3}}}
\newtcolorbox{countrybox}[1]{wbbase, breakable, frame hidden,
  borderline={1.4pt}{0pt}{Ink, line cap=round, dash pattern=on 0pt off 2.4pt},
  title={\icocountry\ #1}}

% ── Padding to a multiple of 4 (copier booklet mode) ────────────────────────
\newcommand{\notespage}{\clearpage\markboth{Notes}{Notes}%
  {\headingfont\bfseries\Large Notes}\par\vspace{3mm}%
  \noindent\begin{tikzpicture}\foreach \y in {0,...,17}{\draw[black!35, line width=0.4pt] (0,-\y*9mm) -- (\linewidth,-\y*9mm);}\end{tikzpicture}}
\newcount\padleft
\newcommand{\padtofour}{\clearpage
  \pgfmathtruncatemacro\pagesused{\the\value{page}-1}%
  \pgfmathtruncatemacro\padn{mod(4-mod(\pagesused,4),4)}%
  \padleft=\padn\relax
  \loop\ifnum\padleft>0 \notespage\advance\padleft by -1 \repeat
  \clearpage}

% ── A5: every question full page width (stacked, never side by side) ───────
\renewcommand{\taskpair}[4]{%
  \par\needspace{6\baselineskip}\vspace{2pt}%
  \task #1\par\vspace{2pt}\answerbox{#2}%
  \par\needspace{6\baselineskip}%
  \task #3\par\vspace{2pt}\answerbox{#4}\par\vspace{2pt}}
\renewcommand{\ptpair}[6]{%
  \pt{#1}#2\par\answerbox{#3}%
  \pt{#4}#5\par\answerbox{#6}}
\renewcommand{\pqpair}[4]{%
  \needspace{4\baselineskip}\vspace{5pt}%
  \stepcounter{pqnum}\noindent\textbf{\thepqnum.}\quad #1\par\vspace{2pt}
  \begin{tcolorbox}[enhanced,arc=2mm,colback=Paper,colframe=Line,left=2mm,right=2mm,
    top=1mm,bottom=1mm,boxrule=1pt,height=#2]\end{tcolorbox}
  \stepcounter{pqnum}\noindent\textbf{\thepqnum.}\quad #3\par\vspace{2pt}
  \begin{tcolorbox}[enhanced,arc=2mm,colback=Paper,colframe=Line,left=2mm,right=2mm,
    top=1mm,bottom=1mm,boxrule=1pt,height=#4]\end{tcolorbox}}

\begin{document}
\pdfbookmark[1]{Cover}{cover}
\setlength{\parindent}{0pt}
\setlength{\parskip}{4pt}

% ---- Cover: the unit's cover art, full bleed, in colour ------------------
\thispagestyle{empty}
\begin{tikzpicture}[remember picture, overlay]
  \node[anchor=north] at (current page.north) {\includegraphics[width=\paperwidth]{../../../covers/science-skills.jpg}};
  \node[anchor=north east, fill=white, draw=Ink, line width=0.6pt, rounded corners=1.2mm,
        inner xsep=2mm, inner ysep=1.8mm, text width=30mm, font=\headingfont\footnotesize]
    at ([xshift=-4mm, yshift=-4mm]current page.north east)
    {\textbf{Name}\par\vspace{4mm}\hrule\vspace{2mm}\textbf{Class}\par\vspace{4mm}\hrule};
  \node[anchor=south east, fill=white, draw=Ink, line width=0.6pt, rounded corners=1.2mm,
        inner xsep=2.6mm, inner ysep=2mm, font=\headingfont\bfseries, align=right]
    at ([xshift=-4mm, yshift=6mm]current page.south east)
    {Advanced Science Skills\\[-1pt]{\normalfont\headingfont\scriptsize Topics 1--3}};
\end{tikzpicture}
\mbox{}
\clearpage

% ---- Inside cover ---------------------------------------------------------
\thispagestyle{plain}
{\headingfont\bfseries\fontsize{22}{26}\selectfont Advanced Science Skills}\par\vspace{1mm}
{\headingfont\large Topics 1--3 \textbullet\ Above the Year 7 level}\par\vspace{1mm}
{\small\itshape VC 2.0: VC2S10I01--I07 \textbullet\ VC2S10H01}\par\vspace{4mm}
\textbf{Teacher}\ \hrulefill\par

\begin{countrybox}{Acknowledgement of Country}
This booklet was created on the unceded Country of the Wurundjeri Woi-wurrung people of the Kulin Nation. We acknowledge Elders past, present and emerging.
\end{countrybox}

\begin{infobox}{How this booklet works}
\textbf{Advanced Science Skills} goes deeper into the ideas of the lesson books: more powerful models, the reasoning behind them, and harder problems that bring several ideas together. Each topic builds on the lessons named under its title. Each opens with its key vocabulary, then:\par\smallskip
\renewcommand{\arraystretch}{1.25}
\begin{tabular}{@{}p{24mm}p{70mm}@{}}
\textbf{Part A} & a reading with comprehension questions \\
\textbf{Part B} & the deeper idea, worked through \\
\textbf{Part C} & practice on your own \\
\textbf{Part D} & a challenge that pulls the ideas together \\
\end{tabular}

\smallskip
\begin{tabular}{@{}p{24mm}p{70mm}@{}}
\textbf{I do} \tagicon{\icodots{1}} & a worked example (solid frame) \\
\textbf{We do} \tagicon{\icodots{2}} & you try one with help (dashed frame) \\
\textbf{Reading} \tagicon{\icoread} & a passage to read (thin frame) \\
\end{tabular}
\end{infobox}

\clearpage
\markboth{Contents}{Contents}
{\headingfont\bfseries\LARGE Contents}\par\vspace{1mm}\noindent\rule{\linewidth}{0.6pt}\par\vspace{2mm}
\renewcommand{\arraystretch}{1.45}
\begin{tabular}{|C{1.1cm}|L{5.2cm}|C{1.7cm}|C{1.4cm}|}
  \hline\thead{Topic} & \thead{Title} & \thead{Deepens lesson} & \thead{Page} \\ \hline
  1 & \hyperlink{lesson.1.1}{How Science Changes Its Mind} & 1 & \hyperlink{lesson.1.1}{p.\ \pageref{sec:1}} \\ \hline
  2 & \hyperlink{lesson.2.1}{Designing Better Experiments} & 4 & \hyperlink{lesson.2.1}{p.\ \pageref{sec:2}} \\ \hline
  3 & \hyperlink{lesson.3.1}{Measuring and Controlling Risk} & 5 & \hyperlink{lesson.3.1}{p.\ \pageref{sec:3}} \\ \hline
  --- & \hyperlink{glossary}{Glossary} & & \hyperlink{glossary}{\pageref{sec:glossary}} \\ \hline
\end{tabular}
\renewcommand{\arraystretch}{1}
\clearpage
\setcounter{section}{0}\renewcommand{\lessonart}{1}
\sectiondiv{How Science Changes Its Mind}{Peer review and replication \textbullet{} Judging a claim \textbullet{} Correlation is not cause}\label{sec:1}
\lessoncodes{Deepens Lesson 1 \textbullet{} VC2S10H01 \textbullet{} VC2S10I06 \textbullet{} VC2S10I07}

\vocabhead
\begin{ctab}
\begin{tabular}{|SW{3.4cm}|W{7.0cm}|}
  \hline\hd & \thead{Word} & \thead{Definition (copy from the slide)} \\ \hline
  & \scaf{\textbf{Peer review}\label{voc:peer-review}} & \scaf{Checking of a scientist's work by other experts before it is published.} \\ \hline
  & \textbf{Replicate}\label{voc:replicate} & \\ \hline
  & \textbf{Correlation}\label{voc:correlation} & \\ \hline
  & \textbf{Cause}\label{voc:cause} & \\ \hline
  & \textbf{Placebo}\label{voc:placebo} & \\ \hline
  & \textbf{Blind trial}\label{voc:blind-trial} & \\ \hline
\end{tabular}
\end{ctab}

\lessonpart{A}{Reading and comprehension}{Teal}\label{read:1}

\begin{readingbox}{Reading --- Too Good to Be True}
In 1989 two chemists announced that they had produced nuclear fusion, the energy source of the Sun, in a jar of water on a laboratory bench. ``Cold fusion'' promised endless cheap energy, and it made headlines around the world. Within weeks, laboratories everywhere tried to \textbf{replicate} the experiment: to repeat it and get the same result. Almost none could. Within a year most scientists had concluded the original results were mistaken.

The same checking works in the other direction. When Barry Marshall and Robin Warren in Perth claimed that most stomach ulcers were caused by bacteria, not stress, other scientists doubted them. Marshall even drank a culture of the bacteria and became ill. As other researchers replicated their findings, the idea was accepted, and in 2005 they won the Nobel Prize.

Before a study is published, other experts check its methods and reasoning: this is \textbf{peer review}. It cannot catch every mistake, but together with replication it is how science corrects itself. One trap is confusing a \textbf{correlation}, two things that tend to happen together, with a \textbf{cause}, where one thing makes the other happen.
\end{readingbox}

\rcreset
\rcq{What does it mean to replicate an experiment?}{2cm}
\rcq{How did replication settle the cold fusion claim?}{2.5cm}
\rcq{Why was the ulcer idea eventually accepted, even though many doubted it at first?}{2.5cm}
\rcq{What is the difference between a correlation and a cause?}{2.5cm}

\lessonpart{B}{Theory and examples}{Blue}\label{theory:1}

\subsection{Judging a Claim}\label{sub:a1.1}

\begin{skillbox}{Five questions to ask of any claim}
\textbf{1. Who says so,} and do they gain from it? \quad \textbf{2. What is the evidence:} a test, or only stories? \quad \textbf{3. How many} were tested? \quad \textbf{4. Was there a comparison group} that did not get the treatment? \quad \textbf{5. Has anyone else replicated it?}
\end{skillbox}

\begin{workedex}{A memory drink}
A company website says its drink ``improves memory by 30\%''. It gave the drink to 10 of its own staff, who all said they felt sharper.

\smallskip
The company profits if you believe it. The evidence is only how people felt, not a measured test. Ten people is a tiny sample. Nobody drank a plain drink for comparison, and no one else has repeated it. The claim is \textbf{weak} on every question.
\end{workedex}

\begin{wedo}{A study app}
A news report says students who use a study app score higher on tests. The study, by a university, compared 2000 users with 2000 non-users. Use the five questions to judge the claim. What else would you want to know?
\answerbox{5cm}
\end{wedo}

\subsection{Correlation Is Not Cause}\label{sub:a1.2}

\begin{workedex}{Ice cream and shark attacks}
On days when more ice cream is sold, there are more shark attacks. Does ice cream cause shark attacks?

\smallskip
No. Both go up on hot days, when more people buy ice cream \emph{and} more people swim. A third factor, the weather, causes both. The two are correlated, but one does not cause the other.
\end{workedex}

\begin{wedo}{Bigger shoes, better spelling}
Across a primary school, students with bigger shoe sizes are better at spelling. Explain why, and name the third factor.
\answerbox{4cm}
\end{wedo}

\lessonpart{C}{Practice}{Amber}\label{prac:1}

\begin{keepext}
\task Most doctors did not accept Marshall and Warren's idea at first. Explain why scientists are right to ask for strong evidence before they change an accepted idea, and why they must still be willing to change it.
\answerbox{5cm}
\end{keepext}
\begin{keepext}
\task A website says that wearing a copper bracelet cures sore joints. Describe how a scientist could test this claim fairly.
\answerbox{5cm}
\end{keepext}


\lessonpart{D}{Challenge}{Violet}\label{ext:1}

\task A company has a new cream it says eases sore muscles. Design a fair trial. Explain how you would use a \textbf{placebo} (a cream that looks and smells the same but has no active ingredient), why the people in the trial should not know which cream they have (a \textbf{blind trial}), and why it is even better if the researchers checking them do not know either.
\answerbox{9cm}

\clearpage

% ══════════════════════════════════════════════════════════════════════════════
%  TOPIC 2 — DESIGNING BETTER EXPERIMENTS                    (deepens Lesson 4)
% ══════════════════════════════════════════════════════════════════════════════
\setcounter{section}{1}\renewcommand{\lessonart}{4}
\sectiondiv{Designing Better Experiments}{Repeats and averages \textbullet{} Outliers \textbullet{} Accuracy and precision \textbullet{} Making results trustworthy}\label{sec:2}
\lessoncodes{Deepens Lesson 4 \textbullet{} VC2S10I02 \textbullet{} VC2S10I03 \textbullet{} VC2S10I05}

\vocabhead
\begin{ctab}
\begin{tabular}{|SW{3.4cm}|W{7.0cm}|}
  \hline\hd & \thead{Word} & \thead{Definition (copy from the slide)} \\ \hline
  & \scaf{\textbf{Trial}\label{voc:trial}} & \scaf{One run of an experiment.} \\ \hline
  & \textbf{Mean}\label{voc:mean} & \\ \hline
  & \textbf{Outlier}\label{voc:outlier} & \\ \hline
  & \textbf{Accurate}\label{voc:accurate} & \\ \hline
  & \textbf{Precise}\label{voc:precise} & \\ \hline
  & \textbf{Reliable}\label{voc:reliable} & \\ \hline
\end{tabular}
\end{ctab}

\lessonpart{A}{Reading and comprehension}{Teal}\label{read:2}

\begin{readingbox}{Reading --- Once Is Never Enough}
Drop a paper helicopter and time its fall: 2.31 seconds. Drop it again: 2.28. Again: 2.35. No two \textbf{trials} give quite the same answer, because no two drops are quite the same: your reaction time, a draught, the angle of release. That is why scientists repeat measurements and take the \textbf{mean}, the average. Random errors in one direction tend to cancel errors in the other.

Sometimes one result is far from the rest: 3.10 seconds, perhaps because the stopwatch was started late. A result that does not fit the pattern is an \textbf{outlier}. Scientists look for a reason, and if it was a mistake, leave it out of the mean and say so.

Two other words matter. A measurement is \textbf{accurate} if it is close to the true value. Measurements are \textbf{precise} if they are close to each other. Kitchen scales that always read 50\,g too heavy are precise but not accurate. Results that come out much the same when the experiment is repeated, by you or someone else, are \textbf{reliable}.
\end{readingbox}

\rcreset
\rcq{Why do repeated trials give slightly different results?}{2.5cm}
\rcq{Why does taking a mean give a better answer than one trial?}{2.5cm}
\rcq{What should a scientist do about an outlier?}{2.5cm}
\rcq{Explain how scales can be precise but not accurate.}{2.5cm}

\lessonpart{B}{Theory and examples}{Blue}\label{theory:2}

\subsection{Means and Outliers}\label{sub:a2.1}

\begin{workedex}{Helicopter times}
Times (s): 2.31, 2.28, 2.35, 3.10, 2.30.

\smallskip
3.10 is far from the others, which are all within 0.07\,s of each other, so it is an \textbf{outlier} (probably a late start). Leave it out:\\
mean $= (2.31 + 2.28 + 2.35 + 2.30) \div 4 = 9.24 \div 4 = \mathbf{2.31\ s}$.
\end{workedex}

\begin{wedo}{Ball bounce heights}
Bounce heights (cm): 42, 44, 29, 43. Identify the outlier, suggest a reason for it, and calculate the mean without it.
\workbox{4.5cm}
\end{wedo}

\subsection{Accurate and Precise}\label{sub:a2.2}

\begin{workedex}{Two thermometers}
Boiling water (100\,\textdegree C) is measured three times with each of two thermometers.
A: 100.1, 99.9, 100.0. \quad B: 96.0, 96.1, 95.9.

\smallskip
A is accurate and precise. B is precise (its readings agree) but not accurate (they are all about 4\,\textdegree C low). B has a \textbf{systematic error}: the same error every time, which repeating cannot fix, but checking against a known value can.
\end{workedex}

\begin{wedo}{Two more thermometers}
C: 98.0, 102.5, 99.5. \quad D: 91.0, 97.0, 94.5. Describe each as accurate or not, and precise or not. Which would you trust least, and why?
\answerbox{4.5cm}
\end{wedo}

\lessonpart{C}{Practice}{Amber}\label{prac:2}

\begin{keepext}
\task Sam tests two paper helicopters. One has longer blades \emph{and} is made from thicker paper. It falls more slowly. Explain why Sam cannot conclude that longer blades make a helicopter fall more slowly, and how to fix the test.
\answerbox{5cm}
\end{keepext}
\begin{keepext}
\task Mia's results did not support a classmate's hypothesis. Explain why a hypothesis that turns out wrong is still useful in science.
\answerbox{5cm}
\end{keepext}


\lessonpart{D}{Challenge}{Violet}\label{ext:2}

\task A student tests whether fertiliser makes bean plants grow taller. She puts one plant with fertiliser on a sunny windowsill and one without on a shady shelf, waters them ``when they look dry'', and measures each once, after a week, with a ruler that starts 5\,mm from the end. She concludes fertiliser makes plants grow taller. Identify four problems with her method and explain how you would fix each.
\answerbox{9cm}

\clearpage

% ══════════════════════════════════════════════════════════════════════════════
%  TOPIC 3 — MEASURING AND CONTROLLING RISK                  (deepens Lesson 5)
% ══════════════════════════════════════════════════════════════════════════════
\setcounter{section}{2}\renewcommand{\lessonart}{5}
\sectiondiv{Measuring and Controlling Risk}{Hazard and risk \textbullet{} The risk matrix \textbullet{} The hierarchy of controls \textbullet{} Writing a risk assessment}\label{sec:3}
\lessoncodes{Deepens Lesson 5 \textbullet{} VC2S10I02 \textbullet{} VC2S10I03}

\vocabhead
\begin{ctab}
\begin{tabular}{|SW{3.4cm}|W{7.0cm}|}
  \hline\hd & \thead{Word} & \thead{Definition (copy from the slide)} \\ \hline
  & \scaf{\textbf{Likelihood}\label{voc:likelihood}} & \scaf{How probable it is that harm will happen.} \\ \hline
  & \textbf{Consequence}\label{voc:consequence} & \\ \hline
  & \textbf{Risk matrix}\label{voc:risk-matrix} & \\ \hline
  & \textbf{Hierarchy of controls}\label{voc:hierarchy-of-controls} & \\ \hline
  & \textbf{Residual risk}\label{voc:residual-risk} & \\ \hline
\end{tabular}
\end{ctab}

\lessonpart{A}{Reading and comprehension}{Teal}\label{read:3}

\begin{readingbox}{Reading --- How Safety Officers Think}
A hazard is anything that could cause harm. Risk is the chance that it actually will, and how bad it would be. Safety officers in laboratories, factories and hospitals judge risk with two questions. What is the \textbf{likelihood} that someone is harmed? And what is the \textbf{consequence}, how serious would the harm be? A \textbf{risk matrix} combines the two answers into a single rating.

Not all ways of reducing risk are equally good. The \textbf{hierarchy of controls} ranks them. Best is to remove the hazard altogether. Next, swap it for something less dangerous. Then use equipment that separates people from the hazard, such as a guard on a machine or a fume cupboard. Below that come rules, signs and training, which only work if people follow them. Last, and least effective, is protective clothing, such as gloves and safety glasses, which only protects the person wearing it, and only if it is worn properly.

After the controls are in place, some risk always remains. This \textbf{residual risk} must be low enough to accept before the work goes ahead.
\end{readingbox}

\rcreset
\rcq{What two questions does a safety officer ask to judge a risk?}{2.5cm}
\rcq{Why is removing a hazard better than wearing gloves?}{2.5cm}
\rcq{Why are rules and signs lower on the hierarchy than a guard on a machine?}{2.5cm}
\rcq{What is residual risk?}{2cm}

\lessonpart{B}{Theory and examples}{Blue}\label{theory:3}

\subsection{The Risk Matrix}\label{sub:a3.1}

\begin{ctab}
\begin{tabular}{|C{2.6cm}|C{2.4cm}|C{2.4cm}|C{2.4cm}|}
  \hline & \thead{Minor (1)} & \thead{Moderate (2)} & \thead{Serious (3)} \\ \hline
  \thead{Likely (3)} & 3 medium & 6 high & 9 high \\ \hline
  \thead{Possible (2)} & 2 low & 4 medium & 6 high \\ \hline
  \thead{Unlikely (1)} & 1 low & 2 low & 3 medium \\ \hline
\end{tabular}
\end{ctab}
\begin{center}\small Risk rating $=$ likelihood $\times$ consequence\end{center}

\begin{workedex}{Heating water with a Bunsen burner}
Hazard: a scald from boiling water. Without controls a spill is \textbf{possible} (2) and a scald is \textbf{moderate} (2): $2 \times 2 = 4$, medium.
With controls (beaker on a tripod and gauze, standing up, hot water left to cool before it is moved), a spill becomes \textbf{unlikely} (1): $1 \times 2 = 2$, low. The residual risk is acceptable.
\end{workedex}

\begin{wedo}{Cutting with a scalpel}
Hazard: cutting a finger while dissecting a flower. Rate the risk before and after controls, and name the controls you would use.
\answerbox{5cm}
\end{wedo}

\subsection{The Hierarchy of Controls}\label{sub:a3.2}

\begin{skillbox}{Most to least effective}
\textbf{1.} Eliminate the hazard \quad \textbf{2.} Substitute something safer \quad \textbf{3.} Engineering controls (guards, fume cupboards) \quad \textbf{4.} Administrative controls (rules, training, signs) \quad \textbf{5.} Personal protective equipment (gloves, glasses, coats)
\end{skillbox}

\begin{workedex}{A strong-smelling solvent}
An experiment uses a solvent whose fumes cause headaches. \textbf{Substitute} a safer solvent if one works; if not, use it in a \textbf{fume cupboard} (engineering), \textbf{train} students and limit the amount (administrative), and wear \textbf{gloves and safety glasses} (PPE). Several controls together give the lowest residual risk.
\end{workedex}

\begin{wedo}{Broken glass}
Students keep cutting themselves on chipped beakers. Suggest one control from each of three different levels of the hierarchy, and say which you think would work best.
\answerbox{5cm}
\end{wedo}

\lessonpart{C}{Practice}{Amber}\label{prac:3}

\begin{keepext}
\task A teacher says: ``We cannot remove every hazard from a lab, but we can make the risk very small.'' Explain what she means, using heating water as an example.
\answerbox{5cm}
\end{keepext}
\begin{keepext}
\task Some chemicals are only used by teachers, in a fume cupboard. Suggest why, and what kind of hazard a fume cupboard controls.
\answerbox{5cm}
\end{keepext}


\lessonpart{D}{Challenge}{Violet}\label{ext:3}

\task Write a full risk assessment for heating salt water to dryness in an evaporating basin over a Bunsen burner. For three hazards, give the likelihood and consequence before controls, the controls (with their level in the hierarchy), and the risk rating after controls.
\begin{ctab}
\begin{tabular}{|W{2.0cm}|C{1.3cm}|W{3.6cm}|C{1.3cm}|}
  \hline\thead{Hazard} & \thead{Risk before} & \thead{Controls (and level)} & \thead{Risk after} \\ \hline
  & & & \\[12mm] \hline
  & & & \\[12mm] \hline
  & & & \\[12mm] \hline
\end{tabular}
\end{ctab}

\clearpage

\clearpage\markboth{Glossary}{Glossary}\hypertarget{glossary}{}\pdfbookmark[1]{Glossary}{glossarybm}
{\headingfont\bfseries\LARGE Glossary}\label{sec:glossary}\par\vspace{1mm}\noindent\rule{\linewidth}{0.6pt}\par
{\small Every key word in this booklet and the page it is introduced on.}\par\smallskip
\renewcommand{\arraystretch}{1.3}
\begin{longtable}{|>{\RaggedRight}p{34mm}|p{9mm}|p{\dimexpr\linewidth-43mm-6\tabcolsep-4\arrayrulewidth-1pt\relax}|}
\hline \rule{0pt}{4mm}\headingfont\bfseries Word & \headingfont\bfseries Page & \headingfont\bfseries In my own words \\ \hline \endhead
  Accurate & \pageref{voc:accurate} & \\ \hline
  Blind trial & \pageref{voc:blind-trial} & \\ \hline
  Cause & \pageref{voc:cause} & \\ \hline
  Consequence & \pageref{voc:consequence} & \\ \hline
  Correlation & \pageref{voc:correlation} & \\ \hline
  Hierarchy of controls & \pageref{voc:hierarchy-of-controls} & \\ \hline
  Likelihood & \pageref{voc:likelihood} & \\ \hline
  Mean & \pageref{voc:mean} & \\ \hline
  Outlier & \pageref{voc:outlier} & \\ \hline
  Peer review & \pageref{voc:peer-review} & \\ \hline
  Placebo & \pageref{voc:placebo} & \\ \hline
  Precise & \pageref{voc:precise} & \\ \hline
  Reliable & \pageref{voc:reliable} & \\ \hline
  Replicate & \pageref{voc:replicate} & \\ \hline
  Residual risk & \pageref{voc:residual-risk} & \\ \hline
  Risk matrix & \pageref{voc:risk-matrix} & \\ \hline
  Trial & \pageref{voc:trial} & \\ \hline
\end{longtable}
\padtofour
\end{document}
